% BHU PYQ -- Manually transcribed & error-corrected
% Paper: BCH-224 : Business Mathematics
% Exam: B.Com. (Hons.) Semester IV Examination, 2022-23

\documentclass[11pt,a4paper]{article}
\usepackage[a4paper,top=2.5cm,bottom=2.5cm,left=2.8cm,right=2.8cm,headheight=14pt]{geometry}
\usepackage{amsmath,amssymb}
\usepackage[T1]{fontenc}
\usepackage[utf8]{inputenc}
\usepackage{lmodern,microtype,fancyhdr,enumitem,hyperref,lastpage,parskip}

\hypersetup{
    colorlinks=true,
    urlcolor=black,
    linkcolor=black,
    pdftitle={BCH-224: Business Mathematics -- BHU B.Com. (Hons.) Sem IV 2022-23}
}

\pagestyle{fancy}
\fancyhf{}
\renewcommand{\headrulewidth}{0.4pt}
\renewcommand{\footrulewidth}{0.4pt}
\fancyhead[L]{\small   \textit{Banaras Hindu University}}
\fancyhead[R]{\small   \textit{BCH-224: Business Mathematics}}
\fancyfoot[C]{\small Page  \thepage\ of \pageref{LastPage}}


\newlist{parts}{enumerate}{2}
\setlist[parts,1]{label=(\alph*),leftmargin=*,itemsep=5pt,topsep=3pt}
\setlist[parts,2]{label=(\roman*),leftmargin=*,itemsep=3pt,topsep=2pt}

\newcommand{\pts}[1]{\hfill   \textit{[#1]}}


\begin{document}


\begin{center}
  {\large\bfseries BANARAS HINDU UNIVERSITY}\\[2pt]
  {
\normalsize B.Com. (Hons.) --- Semester IV Examination, 2022-23}\\[6pt]
  \rule{0.8\linewidth}{0.5pt}\\[6pt]
  {\large\bfseries Paper: BCH-224: Business Mathematics}\\[4pt]
  {
\normalsize Time: 3 Hours\hfill Maximum Marks: 70}
\end{center}

\rule{\linewidth}{0.4pt}

\smallskip



\noindent   \textit{Instructions: Attempt all the questions. All questions carry equal marks (14 marks each).}

\bigskip




\noindent   \textbf{Question 1.}

\begin{parts}
  \item The $p^{ \text{th}}$ and $q^{ \text{th}}$ terms of an A.P. are $c$ and $d$ respectively. Find the $r^{ \text{th}}$ term.\pts{7}
  \item There are $n$ arithmetic means between $1$ and $51$. If the ratio between the $4^{ \text{th}}$ and $7^{ \text{th}}$ means is $3:5$, find $n$.\pts{7}
\end{parts}

\centerline{   \textbf{OR}}


\begin{parts}
  \item Show that the sum of the $5^{ \text{th}}$ term from the beginning and the $5^{ \text{th}}$ term from the end of an A.P. is equal to the sum of the first and last terms.\pts{7}
  \item The sum of $n$ terms of an A.P. is $40$, the common difference is $2$, and the last term is $13$. Find the number of terms.\pts{7}
\end{parts}

\medskip\rule{\linewidth}{0.2pt}




\noindent   \textbf{Question 2.}

\begin{parts}
  \item The $(p + q)^{ \text{th}}$ term of a G.P. is $m$ and the $(p - q)^{ \text{th}}$ term is $n$. Find the $p^{ \text{th}}$ term.\pts{7}
  \item The arithmetic mean between two numbers is $10$ and their geometric mean is $8$. Find the numbers.\pts{7}
\end{parts}

\centerline{   \textbf{OR}}


\begin{parts}
  \item Find the G.P. sequence whose $5^{ \text{th}}$ and $8^{ \text{th}}$ terms are $1$ and $8$.\pts{7}
  \item Find the sum of $n$ terms of the series:\pts{7}
  \[
  7 + 77 + 777 + 7777 + \dots
  \]
\end{parts}

\medskip\rule{\linewidth}{0.2pt}




\noindent   \textbf{Question 3.}

\begin{parts}
  \item Prove that:\pts{7}
  \[
  {}^{n-1}P_r + r \cdot {}^{n-1}P_{r-1} = {}^nP_r
  \]
  \item How many numbers between $100$ and $10,000$ can be formed with the digits $1, 2, 3, 0, 4, 5, 6, 7$?\pts{7}
\end{parts}

\centerline{   \textbf{OR}}


\begin{parts}
  \item How many triangles can be formed by joining the angular points of a decagon? How many diagonals does it have?\pts{7}
  \item If $m = {}^nC_2$, prove that ${}^mC_2 = 3 \cdot {}^{n+1}C_4$.\pts{7}
\end{parts}

\medskip\rule{\linewidth}{0.2pt}




\noindent   \textbf{Question 4.}

\begin{parts}
  \item Find the $9^{ \text{th}}$ term of $(\sqrt{x} - \sqrt{y})^{16}$.\pts{7}
  \item In the expansion of $(1 + x)^{24}$, if the coefficients of the $r^{ \text{th}}$ and $(r + 4)^{ \text{th}}$ terms are equal, find $r$.\pts{7}
\end{parts}

\centerline{   \textbf{OR}}


\begin{parts}
  \item Find the coefficient of $x^{32}$ in the expansion of:\pts{7}
  \[
  \left(x^4 - \frac{1}{x^3}\right)^{15}
  \]
  \item Find the middle term in the expansion of $(3x - y)^8$.\pts{7}
\end{parts}

\medskip\rule{\linewidth}{0.2pt}




\noindent   \textbf{Question 5.}

\begin{parts}
  \item If $f(x) = \log\left(\frac{1+x}{1-x}\right)$, prove that:\pts{5}
  \[
  f(a) + f(b) = f\left(\frac{a+b}{1+ab}\right)
  \]
  \item Evaluate:\pts{4}
  \[
  \lim_{x  o 0} \frac{\sqrt{1+x} - 1}{x}
  \]
  \item Differentiate with respect to $x$:\pts{5}
  \[
  \frac{a^2 + x^2}{\sqrt{a^2 - x^2}}
  \]
\end{parts}

\centerline{   \textbf{OR}}

Evaluate the following integrals with respect to $x$:\pts{5+4+5}

\begin{parts}
  \item $\displaystyle \int \frac{\sqrt{8 + 3\sqrt{x}}}{5x\sqrt{x}} \, dx$
  \item $\displaystyle \int (x\sqrt{x} + 3) \, dx$
  \item $\displaystyle \int x^n \log x \, dx$
\end{parts}

\vfill
\rule{\linewidth}{0.4pt}

\begin{center}\small
  \href{https://bhu-exam.vercel.app/}{Click here to visit our website}\\[4pt]
  \href{https://me-aryan.vercel.app/}{Click here to visit our contributor}\\[4pt]
  \href{https://quashan-me.vercel.app/}{Click here to visit our contributor}
\end{center}

\end{document}
